\chapter*{Chapitre 8 — Efficacité énergétique et RE2020}
\addcontentsline{toc}{chapter}{Chapitre 8 — Efficacité énergétique et RE2020}

% ============================================================
\section*{Énoncés}
% ============================================================

\exercice{Vérification de conformité RE2020 d'un bâtiment de bureaux}

Un bureau de \SI{500}{\m^2} (zone H2b — Nantes) est équipé d'un système de climatisation réversible VRV avec CTA double flux ($\eta_r = 80\,\%$). Les données de la simulation thermique donnent :
\begin{itemize}
  \item Besoin chauffage : $Q_{\text{ch}} = \SI{28\,000}{\kWh \per \year}$
  \item Besoin refroidissement : $Q_{\text{fr}} = \SI{15\,000}{\kWh \per \year}$
  \item Besoin éclairage : $Q_{\text{écl}} = \SI{8\,000}{\kWh \per \year}$
  \item Auxiliaires (pompes, ventilateurs) : $C_{\text{aux}} = \SI{4\,500}{\kWh \per \year}$
\end{itemize}
Données systèmes : SCOP chauffage = 3,8 ; EER refroidissement = 3,2 ; rendements distribution = 0,97.

\begin{enumerate}
  \item Calculer la consommation en énergie finale pour chaque usage.
  \item Calculer le Cep en \si{\kWh_{ep} \per \m^2 \per \year} (fep électricité = 2,3).
  \item Calculer le Cep,nr ($\text{fep,nr}$ électricité VRV = 2,3 avec fraction non renouvelable 0,55 pour le CVC ; fraction non renouvelable éclairage + auxiliaires = 1,0).
  \item Comparer au seuil indicatif RE2020 bureaux (\SI{90}{\kWh_{ep,nr} \per \m^2 \per \year}) et conclure.
\end{enumerate}

\exercice{Comparaison énergétique — simple flux vs double flux}

Un bâtiment résidentiel collectif de \SI{800}{\m^2} est ventilé par extraction simple flux hygro B. On envisage de le remplacer par un système double flux avec récupération ($\eta_r = 85\,\%$). Données :
\begin{itemize}
  \item Débit d'air total : $\dot{V} = \SI{3\,200}{\m^3 \per \hour}$
  \item DJU chauffage (Lyon H1b) : 2\,130 K.j
  \item Température intérieure conventionnelle : \SI{19}{\degreeCelsius}
  \item Coût du kWh chauffage gaz condensation : \SI{0{,}10}{\euro}
\end{itemize}

\begin{enumerate}
  \item Calculer le besoin de chauffage annuel lié à la ventilation en simple flux.
  \item Calculer l'économie de chauffage apportée par le double flux.
  \item En déduire le temps de retour sur investissement si le surcoût du double flux est de \SI{18\,000}{\euro}.
\end{enumerate}

\newpage
% ============================================================
\section*{Corrigés}
% ============================================================

\subsection*{Corrigé — Exercice 1}

\textbf{1. Consommation en énergie finale par usage :}

Chauffage :
\[
  C_{\text{ef,ch}} = \frac{Q_{\text{ch}}}{\text{SCOP} \times \eta_{\text{dist}}} = \frac{28\,000}{3{,}8 \times 0{,}97} = \frac{28\,000}{3{,}686} = \boxed{\SI{7\,596}{\kWh \per \year}}
\]

Refroidissement :
\[
  C_{\text{ef,fr}} = \frac{Q_{\text{fr}}}{\text{EER} \times \eta_{\text{dist}}} = \frac{15\,000}{3{,}2 \times 0{,}97} = \frac{15\,000}{3{,}104} = \boxed{\SI{4\,833}{\kWh \per \year}}
\]

Éclairage (consommation directe) : $C_{\text{ef,écl}} = \boxed{\SI{8\,000}{\kWh \per \year}}$

Auxiliaires : $C_{\text{aux}} = \boxed{\SI{4\,500}{\kWh \per \year}}$

Total énergie finale :
\[
  C_{\text{ef,total}} = 7\,596 + 4\,833 + 8\,000 + 4\,500 = \boxed{\SI{24\,929}{\kWh \per \year}}
\]

\medskip
\textbf{2. Cep (toute l'électricité, fep = 2,3) :}
\[
  \text{Cep} = \frac{C_{\text{ef,total}} \times \text{fep}}{S_{\text{ref}}} = \frac{24\,929 \times 2{,}3}{500} = \frac{57\,337}{500} = \boxed{\SI{114{,}7}{\kWh_{ep} \per \m^2 \per \year}}
\]

\medskip
\textbf{3. Cep,nr :}

En première approche simplifiée (VRV + fraction non renouvelable 0,55 ; éclairage + auxiliaires avec fraction 1,0) :
\[
  \text{Cep,nr} = \frac{(7\,596 + 4\,833) \times 2{,}3 \times 0{,}55 + (8\,000 + 4\,500) \times 2{,}3 \times 1{,}0}{500}
\]
\[
  = \frac{12\,429 \times 1{,}265 + 12\,500 \times 2{,}3}{500} = \frac{15\,722 + 28\,750}{500} = \frac{44\,472}{500} = \boxed{\SI{88{,}9}{\kWh_{ep,nr} \per \m^2 \per \year}}
\]

\medskip
\textbf{4. Conclusion :}

$\text{Cep,nr} = \SI{88{,}9}{\kWh_{ep,nr} \per \m^2 \per \year} < \SI{90}{\kWh_{ep,nr} \per \m^2 \per \year}$ : le bâtiment est \textbf{conforme à la RE2020} sur l'indicateur Cep,nr avec une marge de 1,2~\%. Il conviendra toutefois de vérifier également les indicateurs Bbio, DH et IC avant de valider la conformité globale.

\subsection*{Corrigé — Exercice 2}

\textbf{1. Besoin de chauffage annuel — ventilation simple flux :}

Débit massique d'air :
\[
  \dot{m}_{\text{air}} = \frac{\dot{V}}{3\,600} \times \rho_{\text{air}} = \frac{3\,200}{3\,600} \times 1{,}2 = \SI{1{,}067}{\kg \per \s}
\]

Besoin annuel par la méthode DJU ($\text{DJU} \times 24$ en K.h) :
\[
  Q_{\text{vent,SF}} = \dot{m}_{\text{air}} \times c_p \times (\text{DJU} \times 24) / 1\,000
  = 1{,}067 \times 1\,006 \times (2\,130 \times 24) / 1\,000
\]
\[
  = 1{,}073 \times 51\,120 / 1\,000 \approx \boxed{\SI{54\,852}{\kWh \per \year}}
\]

\textbf{2. Économie par double flux :}
\[
  \Delta Q = \eta_r \times Q_{\text{vent,SF}} = 0{,}85 \times 54\,852 = \SI{46\,624}{\kWh \per \year}
\]

Économie financière (chauffage gaz, $\eta_{\text{syst}} = 0{,}95$) :
\[
  \Delta E = \frac{46\,624}{0{,}95} \times 0{,}10 = \boxed{\SI{4\,908}{\euro \per \year}}
\]

\textbf{3. Temps de retour sur investissement :}
\[
  T_{\text{retour}} = \frac{18\,000}{4\,908} \approx \boxed{3{,}7\text{ ans}}
\]

\begin{methode}
Ce temps de retour est excellent (inférieur à 5 ans), ce qui justifie pleinement l'investissement dans un système double flux avec récupération thermique.
\end{methode}
